% ============================================================
%  NEXIS AI Companion — Methodology
%  Compile with: pdflatex methodology.tex
% ============================================================
\documentclass[12pt,a4paper]{article}

% --- Packages ---
\usepackage[margin=1in]{geometry}
\usepackage{amsmath, amssymb}
\usepackage{graphicx}
\usepackage{booktabs}
\usepackage{array}
\usepackage{longtable}
\usepackage{hyperref}
\usepackage{xcolor}
\usepackage{listings}
\usepackage{caption}
\usepackage{subcaption}
\usepackage{enumitem}
\usepackage{float}
\usepackage{titlesec}
\usepackage{parskip}
\usepackage[T1]{fontenc}
\usepackage{lmodern}

% --- Code listing style ---
\lstset{
  basicstyle=\ttfamily\small,
  keywordstyle=\color{blue}\bfseries,
  commentstyle=\color{gray},
  stringstyle=\color{orange!80!black},
  backgroundcolor=\color{gray!8},
  frame=single,
  rulecolor=\color{gray!40},
  breaklines=true,
  showstringspaces=false,
  captionpos=b,
  numbers=left,
  numberstyle=\tiny\color{gray},
  numbersep=6pt,
  language=Python
}

% --- Hyperlink setup ---
\hypersetup{
  colorlinks=true,
  linkcolor=blue!70!black,
  urlcolor=blue!70!black,
  citecolor=blue!70!black
}

% ============================================================
\begin{document}

% --- Title ---
\begin{center}
  {\LARGE \textbf{NEXIS AI Companion}}\\[6pt]
  {\large \textbf{Methodology}}\\[4pt]
  {\normalsize Multimodal Emotion Recognition for Mental Health Monitoring}\\[2pt]
  \rule{\linewidth}{0.4pt}
\end{center}

\tableofcontents
\newpage

% ============================================================
\section{Introduction and System Overview}
\label{sec:intro}
% ============================================================

NEXIS AI Companion is a mental health support platform that analyses a user's
emotional state from \textit{three simultaneous modalities}—facial video,
speech audio, and written text—and uses the fused result to track mood over
time, generate risk alerts for negative emotional states, and produce
structured weekly wellness reports.

The system is built on a decoupled architecture:

\begin{itemize}[leftmargin=1.5em]
  \item \textbf{Backend:} Python~3.9+, FastAPI, PostgreSQL, and a stack of
        pre-trained HuggingFace Transformer models combined with a custom
        multi-layer perceptron (MLP) meta-classifier.
  \item \textbf{Frontend:} React (Vite), Tailwind~CSS, Chart.js, served on
        port 5173 and communicating with the backend via REST API.
\end{itemize}

The remainder of this methodology is structured as follows:
Section~\ref{sec:architecture} describes the overall system architecture;
Section~\ref{sec:models} details each modality-specific model;
Section~\ref{sec:preprocessing} covers audio and video preprocessing;
Section~\ref{sec:fusion} explains the late-fusion concatenation strategy and
the disgust override gate;
Section~\ref{sec:training} outlines the MLP training configuration and
dataset;
Section~\ref{sec:backend} documents the API design, background task pipeline,
and database schema;
Section~\ref{sec:risk} explains the 14-day distress scoring and risk
composite;
Section~\ref{sec:sentiment} describes text sentiment pre-screening with
VADER; and
Section~\ref{sec:security} covers authentication and data security.

% ============================================================
\section{System Architecture}
\label{sec:architecture}
% ============================================================

\subsection{High-Level Component Diagram}

The end-to-end inference path begins at the user's browser and terminates at
a persistent result in PostgreSQL (Figure~\ref{fig:arch}).

\begin{figure}[H]
\centering
\begin{verbatim}
  User Browser
    |
    | POST /check-in/multimodal
    |   (multipart: video_file + text_input)
    v
  FastAPI  (checkin.py)
    |  1. Save video  ->  backend/uploads/<uuid>.<ext>
    |  2. INSERT MoodEntry(status=uploaded) -> PostgreSQL
    |  3. Return HTTP 202 Accepted  (immediately)
    |  4. Enqueue BackgroundTask
    v
  _run_analysis_in_background()    (own SQLAlchemy thread)
    |
    v
  predict_emotion()                (utils/predict_emotion.py)
    |
    +-- [no text?] -> Whisper ASR   (openai-whisper)
    |
    +-- extract_features()
    |     |-- ffmpeg  -> 16 kHz mono WAV
    |     |-- pydub   -> float32 waveform
    |     |-- Wav2Vec2 (facebook/wav2vec2-base) -> 7-dim audio vector
    |     |-- DistilRoBERTa (j-hartmann/...)    -> 7-dim text vector
    |     |-- MoviePy sample 5-10 frames
    |     |-- ViT (dima806/facial_emotions_...) -> 7-dim video vector
    |     |-- normalize + renormalize
    |     |-- disgust-gate override
    |     `-- concatenate -> 21-dim feature vector
    |
    `-- sklearn MLP -> emotion label + confidence
          |
          v
    UPDATE MoodEntry(emotion, confidence, probabilities, status=analyzed)
    If negative emotion -> INSERT Alert
\end{verbatim}
\caption{End-to-end inference architecture of the NEXIS AI Companion.}
\label{fig:arch}
\end{figure}

\subsection{Decoupled Request Handling}

A key design decision is to return \texttt{HTTP~202~Accepted} immediately
after saving the video and creating the database record. The ML inference
workload (which can take several seconds on CPU) is then executed
asynchronously in a FastAPI \texttt{BackgroundTask}. This ensures that:

\begin{enumerate}[leftmargin=2em]
  \item The ASGI event loop is never blocked, keeping all other API routes
        responsive during inference.
  \item The frontend can optimistically display the entry as
        \textit{Pending Analysis} and poll for a status update.
  \item The background task opens its own SQLAlchemy \texttt{Session} via
        \texttt{SessionLocal()} to guarantee thread safety.
\end{enumerate}

\subsection{MoodEntry State Machine}

Each check-in entry progresses through a three-state lifecycle:

\begin{center}
\texttt{uploaded} $\;\longrightarrow\;$ \texttt{analyzed}\\
\texttt{uploaded} $\;\longrightarrow\;$ \texttt{failed}
\end{center}

If \texttt{predict\_emotion()} raises an unhandled exception, the entry is
marked \texttt{failed} and the traceback is stored in the
\texttt{analysis\_error} column for debugging.

% ============================================================
\section{Modality-Specific Model Architectures}
\label{sec:models}
% ============================================================

The system employs four pre-trained models as \textit{frozen feature
extractors}; none of them are fine-tuned. All four are loaded once at server
startup and kept in memory for the lifetime of the process
(Table~\ref{tab:models}).

\begin{table}[H]
\centering
\caption{Pre-trained models used as frozen feature extractors.}
\label{tab:models}
\begin{tabular}{llll}
\toprule
\textbf{Modality} & \textbf{Model ID (HuggingFace)} & \textbf{Architecture} & \textbf{Output} \\
\midrule
Facial Video & \texttt{dima806/facial\_emotions\_image\_detection} & ViT & 7-class softmax \\
Speech Audio & \texttt{facebook/wav2vec2-base} & Wav2Vec2 & Latent hidden states \\
Text & \texttt{j-hartmann/emotion-english-distilroberta-base} & DistilRoBERTa & 7-class softmax \\
Transcription & \texttt{openai/whisper-base} & Whisper Encoder-Decoder & Transcript string \\
\bottomrule
\end{tabular}
\end{table}

\subsection{Facial Emotion Recognition — Vision Transformer (ViT)}

The \texttt{dima806/facial\_emotions\_image\_detection} model is a
Vision Transformer (ViT) fine-tuned specifically for facial emotion
recognition. It accepts a single \texttt{PIL.Image} object and outputs a
softmax probability distribution over a set of emotion classes. The
\texttt{AutoImageProcessor} handles resizing, normalisation, and patch
tokenisation automatically before feeding the image into the transformer
encoder. The output \texttt{logits} are passed through \texttt{softmax} to
obtain per-class probabilities:

\begin{lstlisting}[caption={ViT inference for a single frame.}]
inputs = processor(images=pil_image, return_tensors="pt").to(device)
with torch.no_grad():
    logits = model(**inputs).logits
probs = torch.nn.functional.softmax(logits, dim=-1).cpu().numpy()[0]
labels = model.config.id2label
return {labels[i]: float(probs[i]) for i in range(len(probs))}
\end{lstlisting}

\subsection{Speech Emotion Feature Extraction — Wav2Vec2}

\texttt{facebook/wav2vec2-base} is a self-supervised speech model trained
with contrastive learning on unlabelled audio. Rather than using a
fine-tuned classification head, NEXIS extracts the raw
\texttt{last\_hidden\_state} tensor and performs \textit{mean and standard
deviation pooling} to obtain a fixed-size embedding:

\begin{lstlisting}[caption={Wav2Vec2 pooling to produce a latent audio vector.}]
inputs = processor(waveform, sampling_rate=16000,
                   return_tensors="pt", padding=True)
with torch.no_grad():
    hidden = model(**inputs).last_hidden_state  # [1, T, 768]

mean = hidden.mean(dim=1)                       # [1, 768]
std  = hidden.std(dim=1)                        # [1, 768]
emb  = torch.cat([mean, std], dim=1)            # [1, 1536]
emb  = emb.squeeze().cpu().numpy()
\end{lstlisting}

The first seven dimensions of this 1536-dimensional embedding are extracted
and treated as proxy ``emotion activations'' for the seven unified labels.
These are raw latent activations—not semantic probabilities—but the MLP
meta-classifier learns to interpret them appropriately during training.

\subsection{Text Emotion Classification — DistilRoBERTa}

\texttt{j-hartmann/emotion-english-distilroberta-base} is a DistilRoBERTa
model fine-tuned on English text for multi-class emotion classification. Its
label vocabulary uses GoEmotions-style names (\textit{joy, sadness, anger,
fear, neutral, surprise, disgust}), which are remapped to the NEXIS unified
label set (see Section~\ref{sec:preprocessing}). The tokeniser handles
input truncation and padding automatically:

\begin{lstlisting}[caption={DistilRoBERTa text classification.}]
inputs = processor(text_input, return_tensors="pt",
                   truncation=True, padding=True).to(device)
with torch.no_grad():
    logits = model(**inputs).logits
probs = torch.nn.functional.softmax(logits, dim=-1).cpu().numpy()[0]
\end{lstlisting}

\subsection{Automatic Speech Recognition — Whisper}

When a check-in includes a video but no textual input, the system
transcribes the audio track using \texttt{openai/whisper-base} (the
multilingual, base-size Whisper Encoder-Decoder model). The audio is first
extracted to a temporary 16~kHz mono WAV file via FFmpeg, then passed to
\texttt{whisper\_model.transcribe()}:

\begin{lstlisting}[caption={Automatic transcription when text input is absent.}]
result = whisper_model.transcribe(temp_wav_path)
transcript = result.get("text", "").strip()
\end{lstlisting}

The resulting transcript is then used as the text modality input for the
DistilRoBERTa classifier.

\subsection{MLP Meta-Classifier}

The meta-classifier is a \textbf{Multi-Layer Perceptron} implemented with
\texttt{sklearn.neural\_network.MLPClassifier} and serialised to
\texttt{metamodels/emotion\_model.pkl} alongside its
\texttt{LabelEncoder} (\texttt{metamodels/emotion\_encoder.pkl}).

\subsubsection{Input}

The classifier receives a flat 21-dimensional feature vector assembled by
concatenating the three modality probability vectors (each 7-dimensional)
in the order \textbf{Video, Audio, Text}:

\begin{equation}
\mathbf{x} = \bigl[\underbrace{v_0, \ldots, v_6}_{\text{Video (7)}},\;
               \underbrace{a_0, \ldots, a_6}_{\text{Audio (7)}},\;
               \underbrace{t_0, \ldots, t_6}_{\text{Text (7)}}\bigr]
\in \mathbb{R}^{21}
\label{eq:featurevec}
\end{equation}

where the index order within each slice follows the unified label set:
\textit{happy, sad, angry, fearful, neutral, surprise, disgust}.

\subsubsection{Output}

The MLP produces a 7-class posterior distribution
$\hat{\mathbf{p}} = [\hat{p}_0, \ldots, \hat{p}_6]$ from which the
dominant emotion label and its confidence are extracted:

\begin{lstlisting}[caption={MLP inference.}]
probs      = model.predict_proba([features])[0]  # shape: (7,)
pred_index = np.argmax(probs)
pred_label = le.inverse_transform([pred_index])[0]
confidence = probs[pred_index] * 100             # percentage
\end{lstlisting}

% ============================================================
\section{Preprocessing Pipelines}
\label{sec:preprocessing}
% ============================================================

\subsection{Audio Extraction — FFmpeg}

FFmpeg is invoked via Python's \texttt{subprocess.run()} using an
\textit{explicit argument list} (no \texttt{shell=True}) to avoid
Windows path-quoting issues. Audio extraction is performed twice: once for
Whisper ASR and once for Wav2Vec2 feature extraction. Both calls use
identical parameters (Table~\ref{tab:ffmpeg}):

\begin{lstlisting}[caption={FFmpeg audio extraction command.}]
subprocess.run([
    "ffmpeg", "-i", str(video_path),
    "-vn",              # strip video stream
    "-acodec", "pcm_s16le",   # uncompressed 16-bit PCM
    "-ar", "16000",     # resample to 16 000 Hz
    "-ac", "1",         # downmix to mono
    temp_audio_path,
    "-y", "-hide_banner", "-loglevel", "error"
], check=True)
\end{lstlisting}

\begin{table}[H]
\centering
\caption{FFmpeg audio extraction parameters.}
\label{tab:ffmpeg}
\begin{tabular}{lll}
\toprule
\textbf{Parameter} & \textbf{Value} & \textbf{Rationale} \\
\midrule
\texttt{-acodec} & \texttt{pcm\_s16le} & Uncompressed PCM required by \texttt{pydub} and Wav2Vec2 \\
\texttt{-ar} & \texttt{16000} Hz & Wav2Vec2 and Whisper are both trained at 16~kHz \\
\texttt{-ac} & \texttt{1} (mono) & Both models require single-channel audio \\
\bottomrule
\end{tabular}
\end{table}

After extraction, \texttt{pydub.AudioSegment} reads the WAV file and converts
the integer sample array to \texttt{numpy.float32}:

\begin{lstlisting}[caption={Conversion from WAV to float32 NumPy array.}]
audio_seg = AudioSegment.from_wav(temp_audio_path)
waveform  = np.array(audio_seg.get_array_of_samples()).astype(np.float32)
os.remove(temp_audio_path)   # clean up immediately
\end{lstlisting}

\subsection{Video Frame Sampling — MoviePy}

\texttt{moviepy.VideoFileClip} is used to sample individual frames at
\textit{adaptively-spaced timestamps} proportional to the clip duration:

\begin{lstlisting}[caption={Adaptive frame sampling with MoviePy.}]
with VideoFileClip(video_path) as clip:
    duration   = clip.duration
    num_frames = min(10, max(5, int(duration)))
    timestamps = np.linspace(0.1, duration - 0.1, num=num_frames)
    frame_preds = [
        get_prediction_probabilities(
            face_m, face_p,
            Image.fromarray(clip.get_frame(t)), "image"
        )
        for t in timestamps
    ]
\end{lstlisting}

\begin{table}[H]
\centering
\caption{Frame-sampling parameters.}
\label{tab:frames}
\begin{tabular}{lll}
\toprule
\textbf{Parameter} & \textbf{Value} & \textbf{Rationale} \\
\midrule
\texttt{num\_frames} & $\max(5,\,\min(10,\,\lfloor\text{duration}\rfloor))$ & Proportional to clip length; at least 5, at most 10 \\
\texttt{timestamps}  & \texttt{linspace(0.1, duration$-$0.1)} & Uniform spacing; 0.1~s offsets avoid blank/cut frames \\
Frame format & \texttt{PIL.Image} from \texttt{get\_frame(t)} & ViT processor requires a PIL image \\
\bottomrule
\end{tabular}
\end{table}

The per-frame emotion probability distributions from valid (non-empty)
predictions are averaged element-wise to produce a single 7-dimensional
video vector:

\begin{equation}
v_{\ell} = \frac{1}{N}\sum_{i=1}^{N} p^{(i)}_{\ell},
\qquad \ell \in \text{UNIFIED\_LABELS}
\end{equation}

where $N$ is the number of valid frame predictions and $p^{(i)}_{\ell}$
is the normalised probability for label $\ell$ from frame $i$.

\subsection{Label Remapping and Renormalisation}
\label{sec:labelnorm}

Each modality uses a different native label vocabulary. A
\texttt{normalize\_probs()} function maps each model's raw output
into the shared seven-class label space
$\mathcal{L} = \{$\textit{happy, sad, angry, fearful, neutral, surprise, disgust}$\}$:

\begin{lstlisting}[caption={Label remapping dictionaries.}]
UNIFIED_LABELS = ['happy','sad','angry','fearful',
                  'neutral','surprise','disgust']

# DistilRoBERTa uses GoEmotions naming convention
text_map = {
    'joy':     'happy',   'sadness':  'sad',
    'anger':   'angry',   'fear':     'fearful',
    'neutral': 'neutral', 'surprise': 'surprise',
    'disgust': 'disgust'
}

# ViT labels already match the unified set (identity map)
video_map = {e: e for e in UNIFIED_LABELS}
\end{lstlisting}

After remapping, each probability distribution is L\textsubscript{1}-renormalised
to ensure it sums to 1:

\begin{equation}
\hat{p}_{\ell} = \frac{p_{\ell}}{\displaystyle\sum_{j \in \mathcal{L}} p_{j}}
\end{equation}

% ============================================================
\section{Multimodal Fusion Strategy}
\label{sec:fusion}
% ============================================================

\subsection{Late Fusion by Concatenation}

NEXIS employs a \textbf{late fusion} paradigm. Each modality is processed
entirely independently through its own pre-trained encoder and normalisation
pipeline. The resulting three 7-dimensional probability vectors are then
\textit{concatenated} to produce the 21-dimensional joint feature vector
$\mathbf{x}$ defined in Equation~\eqref{eq:featurevec}:

\begin{lstlisting}[caption={Late fusion concatenation (the critical line in predict\_emotion.py).}]
return (
    [video[label] for label in UNIFIED_LABELS] +   # dims 0-6
    [audio[label] for label in UNIFIED_LABELS] +   # dims 7-13
    [text[label]  for label in UNIFIED_LABELS]     # dims 14-20
)
\end{lstlisting}

Late fusion is chosen over early or intermediate fusion because:
\begin{enumerate}[leftmargin=2em]
  \item Each modality has a fundamentally different data type (image, waveform,
        token sequence), making early fusion architecturally complex.
  \item The pre-trained encoders are frozen; their internal representations
        are not adapted to each other.
  \item Late fusion retains interpretability: each 7-dimensional slice in
        $\mathbf{x}$ corresponds directly to a specific modality's emotion
        probability distribution.
\end{enumerate}

\subsection{Confidence-Based Disgust Override Gate}

A hard deterministic rule acts as an override gate \textit{before} the MLP
sees the feature vector. If the text modality detects \textit{disgust} with
probability $\geq 0.80$, all three modality vectors are replaced with a
one-hot distribution peaked at \textit{disgust}:

\begin{lstlisting}[caption={Disgust gate override.}]
def confidence_based_override(text_probs):
    return text_probs.get("disgust", 0) >= 0.80

if confidence_based_override(text):
    for lbl in UNIFIED_LABELS:
        val = 1.0 if lbl == 'disgust' else 0.0
        audio[lbl] = val
        video[lbl] = val
        text[lbl]  = val
\end{lstlisting}

This rule handles cases where strong textual evidence of disgust is present
but the facial or audio channels carry ambiguous or contradictory signals.
The threshold of 0.80 was chosen empirically to minimise false activations
while still catching clearly expressed disgust.

% ============================================================
\section{MLP Training Configuration}
\label{sec:training}
% ============================================================

\subsection{Dataset — RAVDESS}

The MLP meta-classifier was trained offline on the
\textbf{RAVDESS} (Ryerson Audio-Visual Database of Emotional Speech and Song)
dataset. RAVDESS contains 2452 audio-visual recordings from 24 professional
North American actors (12 male, 12 female) expressing eight discrete emotions
at two intensity levels.

\begin{table}[H]
\centering
\caption{RAVDESS dataset properties and preprocessing for training.}
\label{tab:ravdess}
\begin{tabular}{ll}
\toprule
\textbf{Property} & \textbf{Details} \\
\midrule
Source dataset & RAVDESS (24 actors, 8 emotions, 2 intensities) \\
Label mapping & 8 original labels $\rightarrow$ 7 unified (\textit{calm} merged into \textit{neutral}) \\
Feature generation & All three frozen encoders run on each clip $\rightarrow$ 21-dim vector \\
Training target & Ground-truth emotion from RAVDESS file-naming convention \\
Reported accuracy & $\approx$89\% on held-out validation split \\
\bottomrule
\end{tabular}
\end{table}

\subsection{Classifier Hyperparameters}

\begin{table}[H]
\centering
\caption{Scikit-learn \texttt{MLPClassifier} configuration.}
\label{tab:mlp}
\begin{tabular}{ll}
\toprule
\textbf{Hyperparameter} & \textbf{Value / Details} \\
\midrule
Model class & \texttt{sklearn.neural\_network.MLPClassifier} \\
Input dimensionality & 21 features \\
Output classes & 7 (happy, sad, angry, fearful, neutral, surprise, disgust) \\
Hidden layer activation & ReLU (default) \\
Output activation & Softmax \\
Loss function & Cross-entropy log-loss (internal to \texttt{MLPClassifier}) \\
Early stopping & Enabled (\texttt{early\_stopping=True}) \\
Validation fraction & 10\% of training data (default) \\
Patience & \texttt{n\_iter\_no\_change}~=~10 consecutive epochs \\
Serialisation & \texttt{joblib.dump()} $\rightarrow$ \texttt{emotion\_model.pkl} ($\approx$323~KB) \\
Label encoding & \texttt{sklearn.preprocessing.LabelEncoder} $\rightarrow$ \texttt{emotion\_encoder.pkl} \\
\bottomrule
\end{tabular}
\end{table}

\subsection{Early Stopping Mechanism}

Scikit-learn's early stopping implementation holds out a fraction of the
training set as an \textit{internal validation set}. Training halts when
the validation score does not improve for \texttt{n\_iter\_no\_change}
consecutive iterations. At termination, the parameter weights from the
epoch with the highest validation score are restored. This is especially
important given the relatively modest size of RAVDESS: early stopping
prevents the MLP from memorising the training set when trained for many
epochs.

\subsection{Device Selection}

All inference is performed on GPU if one is available, falling back to CPU
otherwise:

\begin{lstlisting}[caption={Device selection.}]
device = "cuda" if torch.cuda.is_available() else "cpu"
\end{lstlisting}

The MLP itself (scikit-learn) always uses CPU; only the HuggingFace
Transformer-based encoders benefit from GPU acceleration.

% ============================================================
\section{Backend API and Integration}
\label{sec:backend}
% ============================================================

\subsection{Framework and Structure}

The backend is a \textbf{FastAPI} application with seven registered routers:

\begin{table}[H]
\centering
\caption{Registered API routers.}
\label{tab:routers}
\begin{tabular}{lll}
\toprule
\textbf{Module} & \textbf{Prefix} & \textbf{Responsibility} \\
\midrule
\texttt{auth.py}          & \texttt{/auth}         & Registration, login, JWT issuance, account management \\
\texttt{checkin.py}       & \texttt{/check-in}     & Video upload, background ML analysis, history retrieval \\
\texttt{survey.py}        & \texttt{/survey}       & PHQ-9 mental health questionnaire submission and results \\
\texttt{quick\_thought.py}& \texttt{/quick-thought}& Short journal entries with VADER sentiment scoring \\
\texttt{dashboard.py}     & \texttt{/dashboard}    & Summary, weekly report, 14-day risk composite \\
\texttt{alerts.py}        & \texttt{/alerts}       & Alert listing, acknowledgement, and backfill \\
\texttt{connections.py}   & \texttt{/connections}  & Guardian/doctor account linking \\
\bottomrule
\end{tabular}
\end{table}

\subsection{Database Models (SQLAlchemy ORM)}

All persistent data is stored in a PostgreSQL database managed through
SQLAlchemy ORM. The key models are:

\begin{itemize}[leftmargin=1.5em]
  \item \textbf{User}: Stores credentials (\texttt{password\_hash} via
        bcrypt), email, display name, and role (\texttt{user},
        \texttt{guardian}, \texttt{doctor}).
  \item \textbf{MoodEntry}: Per check-in record containing the video path,
        optional text input, analysis status, predicted emotion, confidence,
        and full probability distribution (JSON column).
  \item \textbf{Alert}: Auto-generated or backfilled record linking a
        negative \texttt{MoodEntry} to a user-visible alert with urgency
        classification (\textit{Medium} or \textit{High}).
  \item \textbf{SurveyResult}: Stores the total PHQ-9 score, per-question
        answers (JSON), and interpretation string.
  \item \textbf{QuickThought}: Short journal entry with a
        floating-point VADER compound sentiment score.
  \item \textbf{WeeklyReport}: Cached LLM-generated insight report (JSON
        payload) scoped to a 14-day period, avoiding repeated Groq API calls.
\end{itemize}

\subsection{Alert Generation and Backfill}

Alerts are created in two places:

\begin{enumerate}[leftmargin=2em]
  \item \textbf{Inline} — immediately inside \texttt{\_run\_analysis\_in\_background()} when the
        predicted emotion is a member of
        $\{\text{sad, angry, fearful, disgust}\}$.
  \item \textbf{Backfill} — on every \texttt{GET /alerts} request, the
        \texttt{\_backfill\_alerts()} function scans all analysed negative
        \texttt{MoodEntry} rows that do not yet have a linked \texttt{Alert}
        row and creates the missing records. This handles legacy entries and
        any race conditions in the background task.
\end{enumerate}

Urgency is assigned as follows:

\begin{equation}
\text{urgency} =
\begin{cases}
\text{High}   & \text{if } e \in \{\text{fearful, angry}\} \\
\text{Medium} & \text{if } e \in \{\text{sad, disgust}\}
\end{cases}
\end{equation}

\subsection{JWT Authentication}

All protected routes require an \texttt{Authorization: Bearer <token>} header.
Tokens are issued at login using \texttt{python-jose} with the
\texttt{HS256} algorithm and a configurable expiry
(\texttt{ACCESS\_TOKEN\_EXPIRE\_MINUTES}). The dependency
\texttt{get\_current\_user()} decodes the token, validates the subject claim,
and returns the corresponding \texttt{User} ORM object—raising
\texttt{HTTP~401} on any failure.

% ============================================================
\section{14-Day Distress Scoring and Risk Composite}
\label{sec:risk}
% ============================================================

\subsection{Per-Entry Distress Score}

For each analysed \texttt{MoodEntry}, the \texttt{compute\_daily\_score()}
function produces a scalar distress score $s \in [0, 100]$:

\begin{equation}
s = 100 \times \bigl(0.7 \cdot p_{\text{neg}} \;+\; 0.3 \cdot (1 - c)\bigr)
\end{equation}

where $p_{\text{neg}} = \sum_{\ell \in \{\text{sad, angry, fearful, disgust}\}} p_\ell$
is the total probability mass assigned to negative emotions, and
$c \in [0, 1]$ is the normalised model confidence. The uncertainty penalty
$(1 - c)$ ensures that low-confidence predictions are treated as riskier, as
the model is less certain about the user's positive state.

\subsection{14-Day Aggregate Statistics}

The \texttt{aggregate\_last\_14\_days()} function computes the following
statistics over all entries in the rolling 14-day window:

\begin{table}[H]
\centering
\caption{14-day aggregate statistics returned by the dashboard.}
\label{tab:agg}
\begin{tabular}{ll}
\toprule
\textbf{Statistic} & \textbf{Definition} \\
\midrule
\texttt{avg\_score}    & Mean distress score across all entries \\
\texttt{std\_dev}      & Standard deviation of distress scores \\
\texttt{best}          & Minimum distress score (lowest distress) \\
\texttt{worst}         & Maximum distress score (highest distress) \\
\texttt{neg\_ratio}    & Fraction of entries with a predominant negative emotion \\
\texttt{trend\_slope}  & $s_{\text{last}} - s_{\text{first}}$ (positive = worsening) \\
\bottomrule
\end{tabular}
\end{table}

\subsection{Weekly Risk Composite}

The dashboard endpoint \texttt{GET /dashboard/weekly-report} computes a
composite risk score by blending the 14-day aggregate with the PHQ-9
questionnaire score (if available within the same 14-day window):

\begin{equation}
\text{risk} =
\begin{cases}
0.5 \cdot \overline{s} + 0.3 \cdot (r \times 100) + 0.2 \cdot s_{\text{PHQ}}
  & \text{if PHQ-9 available} \\[4pt]
0.6 \cdot \overline{s} + 0.4 \cdot (r \times 100)
  & \text{otherwise}
\end{cases}
\end{equation}

where $\overline{s}$ is the 14-day average distress score,
$r$ is the negative ratio, and
$s_{\text{PHQ}} = ({\text{PHQ-9 raw score}} / 27) \times 100$
is the PHQ-9 score normalised to the same 100-point scale.

\subsection{LLM-Generated Narrative Insights (Groq)}

The weekly report is enriched with structured natural-language insights
generated by the Groq API using the \texttt{llama-3.3-70b-versatile} model.
The function \texttt{generate\_structured\_insights()} builds a JSON prompt
containing the 14-day aggregate, the risk score, and the PHQ-9 result, and
requests a structured JSON response with the following keys:
\textit{summary}, \textit{mood\_direction}, \textit{key\_insights},
\textit{suggestions}, \textit{strengths}, \textit{possible\_triggers}, and
\textit{recommend\_followup}. Generated reports are cached as
\texttt{WeeklyReport} rows in PostgreSQL to avoid redundant Groq API calls.

% ============================================================
\section{Quick Thought VADER Sentiment Scoring}
\label{sec:sentiment}
% ============================================================

In addition to the multimodal video check-in, users can submit short
free-text \textit{Quick Thoughts} (journal entries). These are scored
immediately using the \textbf{VADER} (Valence Aware Dictionary and sEntiment
Reasoner) lexicon-based sentiment analyser from the
\texttt{vaderSentiment} library:

\begin{lstlisting}[caption={VADER sentiment scoring for journal entries.}]
from vaderSentiment.vaderSentiment import SentimentIntensityAnalyzer

analyzer = SentimentIntensityAnalyzer()
scores   = analyzer.polarity_scores(text_content)
sentiment_score = scores["compound"]   # range: [-1.0, +1.0]
\end{lstlisting}

The compound score is stored in \texttt{QuickThought.sentiment\_score} and
fed into the dashboard's \textit{Insight Message} generator, which maps the
score to a human-readable qualitative statement about the user's current
emotional state.

VADER was chosen for journal entries over the transformer-based text model
because it is:
\begin{itemize}[leftmargin=1.5em]
  \item \textbf{Rule-based and instant:} No model inference latency; the
        result is available synchronously within the request.
  \item \textbf{Interpretable:} Positive and negative word contributions are
        lexically defined, not latent.
  \item \textbf{Calibrated for social text:} VADER handles punctuation,
        capitalisation, and emoticons well, which is appropriate for
        informal journal-style input.
\end{itemize}

% ============================================================
\section{Security and Compliance}
\label{sec:security}
% ============================================================

\subsection{Password Hashing}

All passwords are hashed with \textbf{bcrypt} using \texttt{passlib}'s
\texttt{CryptContext} before storage. No plaintext passwords are ever written
to the database.

\subsection{JWT Token Security}

\begin{itemize}[leftmargin=1.5em]
  \item Signed with \texttt{HS256} using a strong \texttt{SECRET\_KEY}
        loaded from the \texttt{.env} file.
  \item Carry \texttt{sub} (user email) and \texttt{role} claims scoped
        per request.
  \item Expiry is configurable via \texttt{ACCESS\_TOKEN\_EXPIRE\_MINUTES}.
\end{itemize}

\subsection{CORS Policy}

The FastAPI application permits requests only from
\texttt{http://localhost:5173} and \texttt{http://127.0.0.1:5173}
(the Vite development server). All HTTP methods and headers are permitted
within this origin list, and credentials (cookies) are allowed.

\subsection{Background Processing as a Security Measure}

Because ML inference runs asynchronously after the HTTP response is sent,
the API response time is deterministic (save-file + INSERT). This prevents
timing-based side-channel attacks that could otherwise reveal whether
specific video content triggered a longer inference path.

\subsection{Environment Variable Management}

All secrets (database credentials, JWT secret, Groq API key) are loaded
from a \texttt{.env} file via \texttt{python-dotenv} and accessed through
\texttt{os.getenv()}. The application raises a \texttt{ValueError} at
startup if any critical variable is absent, preventing silent misconfiguration.

% ============================================================
\section{Summary}
\label{sec:summary}
% ============================================================

Table~\ref{tab:summary} provides a consolidated summary of the methodological
choices made throughout the NEXIS AI Companion system.

\begin{longtable}{p{4.5cm} p{9cm}}
\caption{Summary of methodological choices in NEXIS.}
\label{tab:summary} \\
\toprule
\textbf{Design Decision} & \textbf{Choice and Rationale} \\
\midrule
\endfirsthead
\multicolumn{2}{c}{\tablename~\thetable{} (continued)}\\
\toprule
\textbf{Design Decision} & \textbf{Choice and Rationale} \\
\midrule
\endhead
\midrule
\multicolumn{2}{r}{\textit{Continued on next page}}\\
\endfoot
\bottomrule
\endlastfoot
Fusion strategy & Late fusion by concatenation, preserving modality independence and interpretability \\
Facial encoder & Frozen ViT fine-tuned on FER-style images; pixel-level features handled by the pre-trained image processor \\
Audio encoder & Frozen Wav2Vec2; mean+std pooling of hidden states provides a compact speech embedding without a fine-tuned head \\
Text encoder & Frozen DistilRoBERTa; label remapping aligns its GoEmotions vocabulary to the NEXIS unified label set \\
Meta-classifier & scikit-learn MLP with early stopping; compact, CPU-deployable, and achieves $\approx$89\% accuracy on RAVDESS \\
ASR & Whisper-base; automatically transcribes audio when no text is provided \\
Disgust gate & Hard rule at $\text{text disgust} \geq 0.80$; handles cases where facial/audio cues are ambiguous \\
Async inference & FastAPI BackgroundTask returning HTTP 202; keeps the event loop unblocked \\
Distress scoring & Weighted combination of negative probability mass and inverse confidence; PHQ-9 integrated when available \\
Journal sentiment & VADER compound score; synchronous and interpretable for short informal text \\
LLM insights & Groq Llama-3.3-70B; structured JSON prompt guarantees parseable output; reports cached in PostgreSQL \\
Authentication & JWT (HS256) via python-jose + bcrypt password hashing \\
Database & PostgreSQL with SQLAlchemy ORM; all ML outputs stored as JSON columns for flexible downstream access \\
\end{longtable}

\end{document}
